% ICA 2027 Computational Methods -- FULL PAPER example (unofficial template).
% Review copy only: do not insert names, affiliations, acknowledgments,
% grant identifiers or personally identifying links anywhere in the PDF.
\documentclass[12pt,letterpaper,oneside]{article}
\usepackage{ica-cm}
\addbibresource{references.bib}

% Update BOTH fields before submission. The CM full-paper limit is 8,000
% words, EXCLUDING tables, references, and appendices (see the official CFP).
\newcommand{\PaperTitle}{Who Should Code Next? Budget-Aware Collaboration for Reliable Content Analysis}
\newcommand{\MainWordCount}{TO BE UPDATED}

\begin{document}
\ICATitlePage{\PaperTitle}{Full Paper}{\MainWordCount}

\begin{ICAAbstract}
Replace this paragraph with a concise summary of the communication-research
problem, research questions, design, empirical findings, and contributions.
Do not assert results that have not been observed. Check the current ScholarOne
submission form for any abstract-length field requirements.
\end{ICAAbstract}
\ICAKeywords{content analysis, computational methods, language models, human--AI collaboration}
\clearpage

\section{Introduction}
Explain why the problem matters for communication research, identify a precise
methodological or theoretical gap, and summarize the contributions. The
collaborative content-analysis framework SCALE offers one example of related
work \parencite{zhao2025scale}. Replace all instructional text before submission.

\section{Literature Review and Research Questions}
Connect prior work to constructs such as coding reliability, measurement
validity, and the division of labor between human and automated coders.
\subsection{Research Questions}
\noindent\textbf{RQ1.} Under what conditions does an additional AI coder improve coding reliability?

\noindent\textbf{RQ2.} How does resource allocation affect measurement quality and cost?

\section{Methods}
\subsection{Data and Codebook}
Describe corpus provenance, sampling, coding units, codebook, human references,
and relevant ethical approvals or privacy protections.
\subsection{Coding Procedure}
Specify agent prompts, model versions, random seeds, budgets, adjudication,
and where human intervention occurs. Cite code and data where appropriate,
using anonymous links during double-blind review.
\subsection{Evaluation}
Define metrics before the analysis. For illustration, an error-aware objective
could be written as
\begin{equation}
  J(\pi)=\mathbb{E}_{x}\bigl[L(\hat{y}_{\pi}(x),y(x))\bigr]
      +\lambda\,\mathbb{E}_{x}\bigl[C_{\pi}(x)\bigr].
  \label{eq:objective}
\end{equation}
Explain all symbols and the estimand; an equation alone does not establish a
research contribution.

\section{Results}
Report results with uncertainty estimates, relevant robustness analyses,
and statistical comparisons. Table~\ref{tab:example} is a FORMATTING EXAMPLE
only: it contains no study results.

\begin{table}[htbp]
  \caption{Illustrative Layout for Comparing Coding Conditions}
  \label{tab:example}
  \centering
  \begin{tabular}{lcc}
    \toprule
    Condition & Coding quality & Cost per item \\
    \midrule
    Fixed allocation & --- & --- \\
    Adaptive allocation & --- & --- \\
    \bottomrule
  \end{tabular}
  \ICANote{Dashes are placeholders; replace with actual estimates and define units.}
\end{table}

\begin{figure}[htbp]
  \caption{Example Figure Placement Within a Single-Column Manuscript}
  \label{fig:example}
  \centering
  \fbox{\parbox[c][1.2in][c]{0.85\linewidth}{\centering
    Replace with an anonymized figure, for example a cost--quality curve.}}
  \ICANote{Replace this explanatory placeholder with a figure note.}
\end{figure}

\section{Discussion}
Interpret what the findings mean for communication scholarship, including
construct validity and practical workflow implications. Distinguish observed
results from interpretation.
\subsection{Limitations and Future Work}
Discuss task selection, model change over time, correlated coder errors,
selection bias, and human-label uncertainty where relevant.

\section{Conclusion}
Conclude with the substantive methodological or theoretical advance.

\printbibliography[title={References}]

% OPTIONAL: A CM full paper may include up to 10 single-spaced appendix pages,
% OR an anonymous online appendix of unrestricted length. Appendix materials
% must be excluded from the 8,000-word main-paper count.
% \clearpage
% \appendix
% \section{Additional Coding Instructions}
% Add only anonymized supplementary material; check appendix page count.

\end{document}
